\hypertarget{runtime-services-and-introspection}{%
\chapter{Runtime Services and
Introspection}\label{runtime-services-and-introspection}}

Introspection is the ability of a program to examine its own state and
structure at runtime. Python's dynamic nature makes it one of the most
introspective languages, providing deep access to its own interpreter
state and the structure of its code.

\hypertarget{sys-the-interpreter-interface}{%
\subsection{\texorpdfstring{40.1 \texttt{sys}: The Interpreter
Interface}{40.1 sys: The Interpreter Interface}}\label{sys-the-interpreter-interface}}

The \passthrough{\lstinline!sys!} module provides variables and
functions that interact strongly with the interpreter.

\hypertarget{runtime-environment}{%
\subsubsection{1. Runtime Environment}\label{runtime-environment}}

\begin{itemize}
\tightlist
\item
  \passthrough{\lstinline!sys.argv!}: The command-line arguments passed
  to the script.
\item
  \passthrough{\lstinline!sys.path!}: The list of strings that specifies
  the search path for modules.
\item
  \passthrough{\lstinline!sys.modules!}: The dictionary that maps module
  names to modules which have already been loaded.
\end{itemize}

\hypertarget{resource-management}{%
\subsubsection{2. Resource Management}\label{resource-management}}

\begin{itemize}
\tightlist
\item
  \passthrough{\lstinline!sys.getrefcount(obj)!}: Returns the reference
  count of the object (always one higher than expected because of the
  argument to \passthrough{\lstinline!getrefcount!}).
\item
  \passthrough{\lstinline!sys.getsizeof(obj)!}: Returns the size of an
  object in bytes (calls the \passthrough{\lstinline!tp\_basicsize!} and
  \passthrough{\lstinline!tp\_itemsize!} C slots).
\end{itemize}

\hypertarget{low-level-hooks}{%
\subsubsection{3. Low-Level Hooks}\label{low-level-hooks}}

\begin{itemize}
\tightlist
\item
  \passthrough{\lstinline!sys.settrace(func)!}: Sets the system's trace
  function, allowing you to implement debuggers and code coverage tools.
\item
  \passthrough{\lstinline!sys.setprofile(func)!}: Sets the system's
  profile function for performance analysis.
\end{itemize}

\hypertarget{inspect-deep-object-analysis}{%
\subsection{\texorpdfstring{40.2 \texttt{inspect}: Deep Object
Analysis}{40.2 inspect: Deep Object Analysis}}\label{inspect-deep-object-analysis}}

The \passthrough{\lstinline!inspect!} module provides functions for
learning about live objects.

\hypertarget{type-checking-and-members}{%
\subsubsection{1. Type Checking and
Members}\label{type-checking-and-members}}

\begin{itemize}
\tightlist
\item
  \passthrough{\lstinline!inspect.getmembers(obj)!}: Returns all members
  of an object in a list of \passthrough{\lstinline!(name, value)!}
  pairs.
\item
  \passthrough{\lstinline!inspect.isfunction()!},
  \passthrough{\lstinline!inspect.isclass()!}: Reliable ways to check
  object types.
\end{itemize}

\hypertarget{retrieving-source-code}{%
\subsubsection{2. Retrieving Source Code}\label{retrieving-source-code}}

\passthrough{\lstinline!inspect.getsource(obj)!} retrieves the source
code of a function or class by looking up the filename in the object's
code object and reading from the disk.

\hypertarget{signatures-and-parameters}{%
\subsubsection{3. Signatures and
Parameters}\label{signatures-and-parameters}}

\passthrough{\lstinline!inspect.signature(func)!} returns a
\passthrough{\lstinline!Signature!} object. This is more than just a
list of names; it includes default values, type hints, and the ``kind''
of parameter (positional-only, keyword-only, etc.).

\hypertarget{the-stack-frame}{%
\subsubsection{4. The Stack Frame}\label{the-stack-frame}}

\passthrough{\lstinline!inspect.currentframe()!} and
\passthrough{\lstinline!inspect.stack()!} allow you to walk the
execution stack. You can see which function called the current one,
access its local variables, and even modify them (though this is
extremely dangerous and rarely recommended).

\hypertarget{warnings-managing-runtime-diagnostics}{%
\subsection{\texorpdfstring{40.3 \texttt{warnings}: Managing Runtime
Diagnostics}{40.3 warnings: Managing Runtime Diagnostics}}\label{warnings-managing-runtime-diagnostics}}

The \passthrough{\lstinline!warnings!} module is used to issue alerts
about non-fatal issues (e.g., deprecated features). * \textbf{Filters}:
You can control whether warnings are ignored, printed, or turned into
exceptions using \passthrough{\lstinline!warnings.filterwarnings()!} or
the \passthrough{\lstinline!-W!} command-line switch. *
\textbf{Context}: Warnings include the line of code that triggered them,
making them more useful than simple \passthrough{\lstinline!print()!}
statements for developers.

\hypertarget{ast-programmatic-source-analysis}{%
\subsection{\texorpdfstring{40.4 \texttt{ast}: Programmatic Source
Analysis}{40.4 ast: Programmatic Source Analysis}}\label{ast-programmatic-source-analysis}}

The \passthrough{\lstinline!ast!} module (briefly touched upon in
Chapter 31) allows you to manipulate Python code as a tree structure. *
\textbf{\passthrough{\lstinline!ast.NodeVisitor!}}: A class that you
subclass to traverse the tree and perform actions at specific nodes
(e.g., finding all function calls). *
\textbf{\passthrough{\lstinline!ast.NodeTransformer!}}: A subclass that
allows you to modify the tree, effectively performing
``source-to-source'' compilation or code instrumentation.

\begin{center}\rule{0.5\linewidth}{0.5pt}\end{center}

\begin{center}\rule{0.5\linewidth}{0.5pt}\end{center}
